\section{Related Work}

\subsection{Medication Management as Shared Care Work}

Medication adherence is not only about remembering a dose; it is also about coordinating routines, information, and responsibility. Reminder apps can help people notice a dose. Yet studies of older adults and people managing several chronic conditions show that adherence also depends on treatment complexity, daily habits, and how well technology fits existing routines \cite{27,35,107,108}. Medication work can also include interpreting bodily changes, managing health information, and coordinating treatment across people and settings \cite{8,9,54,85,119}. Conversational systems and studies of family health routines extend this view by showing that self-management may involve relatives and care partners rather than one person acting alone \cite{74,92,102,120}. This makes medication management a useful setting for studying how care work is shared.

Shared medication work is easier to understand when care is treated as an ongoing practice rather than a sequence of isolated choices. Practice-oriented HCI describes technology use as part of everyday activity, while care theory emphasizes repeated adjustment to changing needs \cite{61,80,112}. Research on invisible work likewise shows that patients and caregivers spend substantial effort organizing information, monitoring changes, and keeping care moving \cite{8,47,85,106}. Feminist and relational approaches add that care is shaped by responsibility, responsiveness, and relationships, not only by individual choice \cite{10,43,44,72}. Once medication management is viewed this way, the next question is who gets to decide how that work is divided.

Older adults can rely on others without giving up agency. HCI research challenges accounts that treat older adults as passive recipients, reluctant users, or people whose needs can be inferred from age alone \cite{11,41,48,58,59,64,66}. Interdependence research instead shows that people may depend on others while retaining meaningful control over how help is arranged \cite{13,30,78,101}. Studies of technology adoption and home-based care similarly describe agency as shaped through relationships, social position, and information practices \cite{56,86}. The issue is therefore not whether help exists, but whether the older adult still has a meaningful role in shaping it.

Family care makes this tension visible because support and control can move together. Care-network research has long shown that older adults, relatives, and care workers coordinate information and responsibilities across people \cite{22,63,82,94,118}. More recent work on dementia care, home monitoring, and medication support shows that family involvement can improve coordination while also creating disagreements over access, monitoring, and decision-making \cite{24,36,73,76,81,115}. These studies establish that family support can be valuable without assuming that older adults and caregivers always want the same thing.

\textit{What remains unclear is what happens when the system itself proposes a change in this division of work. Prior research explains reminders, family coordination, and relational agency, but gives less attention to how older adults and caregivers respond when an autonomous agent decides that another person may need to become involved.}

\subsection{Family Visibility, Privacy, and Dignity in Later-Life Care}

Family visibility can support care, but it also changes who knows, who acts, and who feels responsible. Awareness displays and connected care tools can help relatives follow an older adult's daily life \cite{22,82,94,118}. Yet monitoring and information-sharing studies show that the same visibility can create conflict over privacy, data ownership, and expectations to intervene \cite{23,24,65,87,111,115}. Research on shared-device privacy reaches a similar conclusion: access rules are often negotiated between people rather than set once by one user \cite{50,122}. This shifts the problem from whether information should be shared to how that sharing changes relationships.

Those relational changes also affect dignity. Older adults may reject technologies that signal decline, and visible assistive devices can affect social standing or the desire to preserve face \cite{17,20,69,73}. Work on domestic robots similarly shows that older adults weigh dignity, autonomy, and the kind of relationship a system appears to create \cite{21}. Goffman's account of self-presentation helps explain why the social meaning of assistance can matter alongside its practical benefit \cite{34}. A system can therefore preserve formal choice and still feel diminishing if it repeatedly presents a person as dependent.

Because these effects are relational, caregiver perspectives matter without replacing the older adult's account. Caregivers often move among several roles and perform work that systems do not fully capture \cite{47}. Studies of home care and security decisions show that caregivers and care recipients can describe the same arrangement differently, including situations where apparently shared decisions leave the older adult with little practical choice \cite{24,56,65,76}. Family information-sharing can likewise improve coordination while creating new privacy concerns for the person receiving care \cite{115}. These differences lead directly to the question of whose interests a system should prioritize when family members disagree.

Autonomy, privacy, dignity, and caregiver coordination each explain part of that problem, but none explains all of it. Autonomy concerns whether the older adult retains meaningful control; privacy concerns what becomes visible and to whom; dignity concerns how support affects social standing; caregiver research explains how responsibilities are distributed. These concepts help identify the stakes. They do not by themselves specify whose interests an autonomous agent should prioritize at a particular moment.

\textit{The gap is therefore not a lack of theories about family care. It is that existing theories do not fully explain how an agent should move between older-adult control and family involvement, or how people judge whether that shift is appropriate in a particular situation.}

\subsection{Family-Mediated Technology Use in Bangladesh}

In Bangladesh, family involvement is not an added layer around care; it is often part of how care is organized. Research on ageing in Bangladesh and South Asia describes limited formal eldercare, strong reliance on relatives, and changing intergenerational arrangements \cite{5,12,37,51,77,90}. An asset-based perspective cautions against describing this only as a lack of formal services, because households and communities already hold practices through which support is provided \cite{60}. This means that an agent entering the household also enters an existing network of family responsibility.

Technology use in that network is often shared as well. Studies in Bangladesh document collaborative phone use, privacy negotiation on shared devices, and privacy risks that arise through ordinary repair and household practices \cite{1,2,3,4}. Research on intermediated technology use shows that one person may operate, explain, translate, or manage a digital resource for another \cite{33,98}. Such support can be a practical way of gaining access rather than evidence that the person receiving help has no agency. Work with rural women in Bangladesh further shows that technology can both reinforce unequal relations and provide ways to exercise agency within them \cite{110}. Shared use therefore makes authority a household issue, not only an interface issue.

Global South HCI explains why these household arrangements should not be measured against an individual-user model by default. Postcolonial computing and critiques of universal design assumptions show how history, infrastructure, and power shape technology use \cite{26,49}. Structural approaches to digital care similarly show that technology can redistribute existing care work rather than simply add a neutral service \cite{53}. Research on sustainable ICTD in Bangladesh, explainable AI in Global South settings, and Global South perspectives on AI governance further shows that systems are interpreted through local institutions, workflows, and stakeholder positions \cite{84,88,96}. Studies of rural clinical AI and voice assistants add that technical adoption does not guarantee a good fit with everyday practice \cite{89,114}. These findings make the household context central to any account of agent authority.

Existing Bangladesh and ICTD work has therefore established shared devices, mediated use, relational privacy, and household power as important design concerns \cite{1,2,3,4,33,98,110}. Across much of this work, however, the technology remains a tool, information system, or shared device rather than an agent that decides when another person should become involved \cite{26,49,53,84,88,96,114}. That difference matters because initiative changes the question. Families must negotiate not only who can use the system, but whose interests should guide what it does.

\textit{What remains underexplored is how older adults and caregivers in such households negotiate authority when an autonomous agent begins acting within existing family care. Our study examines this gap in Bangladesh without assuming that receiving family help means less agency, or that family-based care means everyone will agree.}

\subsection{From Agent Initiative to Allegiance}

An agent becomes part of the care relationship when it can act before anyone explicitly asks it to. Mixed-initiative research established that systems may balance user control with computational initiative, while newer proactive agents can suggest actions, communicate on a user's behalf, or begin support without a direct command \cite{45,52,70}. Research on plan disclosure and AI decision power shows that people care about when an agent acts, how much authority it has, and whether its intended action is visible beforehand \cite{42,55}. Trust research similarly argues that appropriate reliance depends on understanding what a system can and cannot be trusted to do \cite{57,67}. Initiative therefore raises a question of authority, not only capability.

Contestability and oversight research explains how that authority can remain open to challenge. People need ways to question automated actions, intervene, and understand important system boundaries and decision points \cite{6,7,25,28}. Responsibility research further shows that autonomous action can make it harder to determine who should answer for a decision or its consequences \cite{99}. In family care, this problem becomes sharper because involving a caregiver can change privacy and responsibility at the same time. The next question is whose interests should guide that action.

Alignment research provides one answer by asking what an agent should serve. Foundational work distinguishes instructions, preferences, interests, and values, while formal models show that obedience and benefit do not always point in the same direction \cite{32,38,39,79}. More recent pluralistic approaches recognize that agents may face several legitimate values, communities, or principals rather than one user with one stable preference \cite{29,62,105,121}. This work makes the question ``for whom?'' explicit. Much of it, however, remains formal, conceptual, or model-level rather than an account of how related people negotiate an agent's priority during everyday care.

Care technology reaches the same problem from the opposite direction. Multi-stakeholder systems already involve older adults, relatives, professionals, and caregivers, but many focus on awareness, monitoring, information sharing, or coordination \cite{22,24,63,81,82,94,111,115,118}. Conversational systems can also support medication management or communication between older adults and care partners or providers \cite{74,117}. These systems show that several people can have legitimate interests in the same care process. What they offer less evidence about is how a household responds when the agent itself proposes that responsibility move from one person to another.

We use \emph{allegiance} to name that specific interaction problem: whose interests and authority the agent prioritizes at a particular moment. Allegiance is narrower than alignment and different from a permission setting. Autonomy asks whether the older adult retains control; privacy asks what others can see; consent asks whether an action is allowed; alignment asks which instructions, values, or principals guide the agent. Allegiance focuses on how those concerns meet when more than one person has a legitimate claim over what the agent should do. An older adult may want private support in one situation, shared support in another, and no family involvement in a third. The priority can change.

\textit{The central gap is empirical: plural-alignment research identifies the problem of multiple principals, while family-care research documents multiple stakeholders. We still know little about how older adults and caregivers negotiate an agent's changing priority between them during everyday use, or what makes such a shift accepted, resisted, or reversed.}

\subsection{Making Allegiance Visible, Revisable, and Negotiable}

A changing allegiance is only meaningful if people can notice and question the change. Contestable-AI research argues that automated decisions should remain open to human challenge, while work on explanations and oversight shows the value of making plans, limits, and intervention points clear \cite{6,7,25,28,42,52}. For a shared-care agent, this means more than explaining why a reminder appeared. People also need to know when the agent is about to involve someone else and what that involvement will change.

That visibility must be paired with the ability to change one's mind. Research on distributed and interactive consent argues that decisions involving several people require permission that can be discussed, updated, and withdrawn \cite{71,83,104,109,122}. Studies with older adults similarly show that privacy boundaries change with circumstances rather than remaining fixed after setup \cite{50}. Voice interfaces add another difficulty because a short verbal agreement can become routine without remaining meaningful \cite{100}. Rather than treating caregiver access as permanent, shared-care agents therefore need ways for people to reconsider who is involved and under what conditions.

Making role changes visible can also create social costs. Research on stigma, dignity, and self-presentation shows that visible assistance can signal dependence or affect social standing \cite{17,20,21,34,69}. Invisible-work research adds that making care work visible can change the expectations surrounding that work \cite{85,106}. Relational and feminist ethics therefore caution against treating transparency as automatically beneficial \cite{10,43,44}. The design challenge is to make a change clear enough to question without repeatedly presenting the older adult as incapable.

Even small interface cues can shape these relationships. Research on nudging and persuasive design shows that prompts take meaning from the practices in which people receive them \cite{18,31}. Gamification research likewise cautions that points and streaks are not neutral mechanics, while studies of communication and family fitness show that shared metrics can become signals of closeness, obligation, encouragement, or loss \cite{19,40,46,68,91,97}. Self-determination research connects interactive systems to autonomy, competence, relatedness, and motivation, while warning against using these concepts as generic explanations for engagement \cite{14,95,113,116}. The broader lesson is that an agent's signals about responsibility may shape family relationships as well as individual behavior.

These relational shifts also affect how the problem should be studied. Research on reflexive thematic analysis stresses that qualitative interpretation should make analytic choices explicit rather than treating coder agreement as the only sign of quality \cite{15,16,93}. Work on positionality and reflexivity in HCI similarly argues that researchers should account for how their standpoint shapes interpretation \cite{75,103}. This matters when older adults and caregivers may give different accounts of the same care episode. Keeping those accounts distinct allows disagreement to remain evidence rather than being averaged into one household view.

\textit{What remains is a design and methodological gap. We need to understand how an agent's changing allegiance can be made clear, open to withdrawal or reversal, and negotiable without undermining dignity. We also need to study those shifts through the separate accounts of the people affected by them. This is the gap addressed by our caregiver-inclusive household study and by our analysis of the agent as tool, coach, or advocate.}
